\documentclass[letterpaper,11pt]{moderncv}
% \usepackage[hidelinks]{hyperref}
\usepackage[utf8]{inputenc}
% \usepackage{moderncv}
\usepackage[left=0.7in, right=0.75in, top=1.0in, bottom=1.0in]{geometry}
\usepackage{multicol}
\usepackage{xcolor}
\usepackage{color}


% ModernCV themes
\moderncvstyle{fancy} % Style options are 'casual' (default), 'classic', 'oldstyle', and 'banking'
\moderncvcolor{blue} % Color options 'blue' (default), 'orange', 'green', 'red', 'purple', 'grey', and 'black'

% Personal information
\name{Yang-yang}{Tan}
\address{Born 29 Jan 1997, Hefei, China\\
No.2 Linggong Road, Ganjingzi District, Dalian City, Liaoning Province, P.R.C., 116024}
\email{yangyang-tan@foxmail.com}
\phone[mobile]{(86)18342209877}

\begin{document}

\makecvtitle
%-----------NOTES-----------
\section{Notes}
\cvitem{}{Applying for: Special Postdoctoral Researcher, RIKEN iTHEMS.}

%-----------EDUCATION-----------
\section{Education}
\cventry{September 2019 -- July 2025}{Ph.D., Theoretical Physics}{Dalian University of Technology}{Dalian, Liaoning Province, China}{}{Thesis: Towards Real-Time Dynamics of Critical Phenomena in
 Strong Correlated Systems Advisor: Prof. Wei-jie Fu}
\cventry{September 2015 -- July 2019}{B.S., Physics}{Dalian University of Technology}{Dalian, Liaoning Province, China}{}{ Thesis: The Deep Learning Method in Gluon Spectra Function \phantom{Spectra Function Spectra} Advisor: Prof. Wei-jie Fu}

%-----------WORK EXPERIENCE-----------
\section{Professional Experience}
\cventry{April 2026 -- Present}{Postdoctoral Researcher}{The University of Tokyo}{Tokyo, Japan}{}{}
\cventry{December 2025 -- March 2026}{Visiting Scholar}{Tsinghua University}{Beijing, China}{}{}

%-----------RESEARCH INTERESTS-----------
\section{Research Interests}
\cvlistitem{\textbf{Machine Learning for Physics:} Deep learning for solving inverse problems, Deep learning for quantum field theory, diffusion models}
\cvlistitem{\textbf{Real-time Quantum Field Theory:} Real-time dynamics of quantum field theory, non-perturbative method, dynamical critical phenomena, Non-equilibrium evolution of quantum fields.}
\cvlistitem{\textbf{Phase Transitions and Critical Phenomena:} Critical exponents, dynamical critical exponents, numerical methods for solving fixed-point equations.}
\cvlistitem{\textbf{Relativistic Heavy Ion Collision Phenomenology:} Non-perturbative aspect of strong interacting systems, QCD phase diagram and phase structure.}
%-----------ACADEMIC COLLABORATION-----------
\section{Academic Collaboration}
\cvitem{}{\textbf{DM-LFT Collaboration} (Diffusion Model-Lattice Field Theory): G. Aarts, S. Chen, D. E. Habibi, A. Ipp, R. Kapust, B. Lucini, D. M\"uller, J. M. Pawlowski, T. R. Ranner, Y.-y. Tan, L. Wang, W. Wang, K. Zhou, and Q. Zhu. (2026)}
\cvitem{}{\href{https://fqcd-collaboration.github.io}{\textcolor{blue}{\textbf{fQCD Collaboration}}} (Functional QCD): J. Braun, Y.-r. Chen, W.-j. Fu, F. Gao, C. Huang, F. Ihssen,  K. Kockler, Y. Lu, J. M. Pawlowski, F. Rennecke, F. Sattler, J. Stoll, Y.-y. Tan, Z.-n. Wang, R. Wen, J. Wessely, S. Yin, H.-w. Zheng and N. Zorbach. (2026)}


%-----------PUBLICATIONS-----------
\section{Publications}
\cvitem{}{INSPIRE Profile: \href{https://inspirehep.net/authors/1889875}{\textcolor{blue}{https://inspirehep.net/authors/1889875}}}

\cvitem{15}{\textbf{Y.-y. Tan*}, W.-j. Fu*, L. He*, L. Wang*, (2026) \textit{Solving Functional Renormalization Group Equations with Neural Networks} [\href{https://arxiv.org/abs/2603.21151}{\textcolor{blue}{arXiv:2603.21151}}]}

\cvitem{14}{\textbf{Y.-y. Tan}, S. Yin, Y.-r. Chen,C. Huang, W.-j. Fu*, (2025) \textit{Real-time evolution of critical modes in the QCD phase diagram} [\href{https://arxiv.org/abs/2512.03614}{\textcolor{blue}{arXiv:2512.03614}}]}

\cvitem{13}{Y.-r. Chen*, W.-j. Fu*, \textbf{Y.-y. Tan*}, (2025) \textit{Density functional theory of renormalization group in nuclear matter} [\href{https://arxiv.org/abs/2508.02575}{\textcolor{blue}{arXiv:2508.02575}}]}

\cvitem{12}{W.-j. Fu, C. Huang*, J. M. Pawlowski, \textbf{Y.-y. Tan}, L.-j. Zhou, (2025) \textit{Four-quark scatterings in QCD III}\href{{https://doi.org/10.1103/4sh5-w4yc}}{, \textcolor{blue}{Phys.Rev.D 112 (2025) 5, 054047}}  [\href{https://arxiv.org/abs/2502.14388}{\textcolor{blue}{arXiv:2502.14388}}]}

\cvitem{11}{Y.-r. Chen*, \textbf{Y.-y. Tan}, W.-j. Fu, (2024) \textit{Critical dynamics of Model H within the real-time fRG approach}\href{{https://doi.org/10.1103/PhysRevD.111.094025}}{, \textcolor{blue}{Phys.Rev.D 111 (2025) 9, 094025}} [\href{https://arxiv.org/abs/2406.00679}{\textcolor{blue}{arXiv:2406.00679}}]}

\cvitem{10}{\textbf{Y.-y. Tan}, Y.-r. Chen, W.-j. Fu*, W.-j. Li*, (2024) \textit{Universality of pseudo-Goldstone damping near critical points},  \href{https://doi.org/10.1038/s41467-025-58170-1}{\textcolor{blue}{Nat Commun. 16, 2916 (2025)}} [\href{https://arxiv.org/abs/2403.03503}{\textcolor{blue}{arXiv:2403.03503}}]}



\cvitem{9}{W.-j. Fu, C. Huang*, J. M. Pawlowski, \textbf{Y.-y. Tan}, (2024) \textit{Four-quark scatterings in QCD II}, \href{{https://doi.org/10.21468/SciPostPhys.17.5.148}}{\textcolor{blue}{SciPost Phys. 17, 148 (2024)} }[\href{https://arxiv.org/abs/2401.07638}{\textcolor{blue}{arXiv:2401.07638}}]}

\cvitem{8}{Y.-r. Chen,\textbf{Y.-y. Tan}, W.-j. Fu*, (2023) \textit{Critical dynamics within the real-time fRG approach}\href{{https://doi.org/10.1103/PhysRevD.109.094044}}{, \textcolor{blue}{Phys.Rev.D 109 (2024) 9, 094044}} [\href{https://arxiv.org/abs/2312.05870}{\textcolor{blue}{arXiv:2312.05870}}]}


\cvitem{7}{J. Braun, Y.-r. Chen, W.-j. Fu, F. Gao, F. Ihssen, C. Huang, J. M. Pawlowski, F. Rennecke*, F. Sattler, \textbf{Y.-y. Tan}, R. Wen, S. Yin, (2023) \textit{Soft modes in hot QCD matter}\href{{https://doi.org/10.1103/PhysRevD.111.094010}}{, \textcolor{blue}{Phys.Rev.D 111 (2025) 9, 094010}} [\href{https://arxiv.org/abs/2310.19853}{\textcolor{blue}{arXiv:2310.19853}}]}

\cvitem{6}{C. Huang, \textbf{Y.-y. Tan}, R. Wen, S. Yin, W.-j. Fu*, (2023) \textit{Reconstruction of baryon number distributions}\href{{https://doi.org/10.1088/1674-1137/aceee1}}{, \textcolor{blue}{Chin. Phys. C. 47 (2023) 10, 104106}} [\href{https://arxiv.org/abs/2303.10869}{\textcolor{blue}{arXiv:2303.10869}}]}


\cvitem{5}{S. Yin, \textbf{Y.-y. Tan}, W.-j. Fu*, (2023) \textit{Critical phenomena and functional renormalization group}\href{{https://dx.doi.org/10.11889/j.0253-3219.2023.hjs.46.040002}}{, \textcolor{blue}{NUCLEAR TECHNIQUES. 2023,46(04):040002-040002.}} [\href{{http://chinaxiv.org/abs/202306.00144}}{\textcolor{blue}{ChinaXiv:202306.00144}}](Chinese Review)}

\cvitem{4}{{\textbf{Y.-y. Tan}, C. Huang, Y.-r. Chen, W.-j. Fu*}, $\mathrm{(2022)}$ \textit{Criticality of the $O(N)$ universality via global solutions to nonperturbative fixed-point equations
    }  \href{{https://doi.org/10.1140/epjc/s10052-024-13291-7}}{\textcolor{blue}{Eur. Phys. J. C 84, 897 (2024)}}[\href{{https://arxiv.org/abs/2211.10249}}{\textcolor{blue}{arXiv:2211.10249}}]}
    
    
\cvitem{3}{{W.-j. Fu, C. Huang*, J. M. Pawlowski, \textbf{Y.-y. Tan}}, $\mathrm{(2022)}$ \textit{Four-quark scatterings in QCD I}{,  \href{{https://doi.org/10.21468/SciPostPhys.14.4.069}}{\textcolor{blue}{SciPost Phys. 14, 069 (2023)} }[\href{{https://arxiv.org/abs/2209.13120}}{\textcolor{blue}{arXiv:2209.13120}}]}}
\cvitem{2}{{J. Braun, Y.-r. Chen, W.-j. Fu, A. Gei\ss{}el, J. Horak, C. Huang, F. Ihssen, J. M. Pawlowski, M. Reichert, F. Rennecke, \textbf{Y.-y. Tan}, S. T\"opfel, J.Wessely, N. Wink*, $\mathrm{(2022)}$} \textit{Renormalised spectral flows}\href{{https://doi.org/10.21468/SciPostPhysCore.6.3.061}}{, \textcolor{blue}{SciPost Phys.Core 6 (2023) 061}} [\href{{https://arxiv.org/abs/2206.10232}}{\textcolor{blue}{arXiv:2206.10232}}]}
    
    
\cvitem{1}{\textbf{Y.-y. Tan}, Y.-r. Chen, W.-j. Fu* $\mathrm{(2021)}$ \textit{Real-time dynamics of the $O(4)$ scalar theory within the fRG approach}{,  \href{{https://doi.org/10.21468/SciPostPhys.12.1.026}}{\textcolor{blue}{SciPost Phys. 12, 026 (2022)} }[\href{{https://arxiv.org/abs/2107.06482}}{\textcolor{blue}{arXiv:2107.06482}}]}}

\section{Contributed Talks}
\cvitem{1}{\textit{Universal Pseudo-Goldstone Damping from the Real-Time fRG}\newline{}\href{https://indico.ihep.ac.cn/event/25688/}{\textcolor{blue}{The QCD Phase Diagram: From Theory to Experimental Signatures}}
\newline{} 8 to 11 October 2025}

\cvitem{2}{\textit{Universal Pseudo-Goldstone Damping from the Real-Time Functional Renormalization Group}\newline{}\href{https://ithems.riken.jp/en/events/universal-pseudo-goldstone-damping-from-the-real-time-functional-renormalization-group}{\textcolor{blue}{RIKEN Seminar}}
\newline{} 26 June 2025 (invited talk)}


\newpage


\cvitem{3}{\textit{Can we infer probability distributions from cumulants? Probabilistic approaches to inverse problems}\newline{}
\href{https://ithems.riken.jp/en/events/can-we-infer-probability-distributions-from-cumulants-probabilistic-approaches-to-inverse-problems}{\textcolor{blue}{DEEP-IN Seminar at RIKEN}}
\newline{} 18 March 2025 (invited talk)}


\cvitem{4}{\textit{Real-time dynamics of pseudo-Goldstone and critical mode in QCD}\newline{}\href{https://indico.cern.ch/event/1374713/}{\textcolor{blue}{The 20th International Conference on QCD in Extreme Conditions(XQCD 2024)}}
\newline{} 17 to 19 July 2024}



\cvitem{5}{\textit{Methods and Applications of Solving Inverse Problems}\newline{} \href{https://snst.lzu.edu.cn/xueshujiaoliu/2023/1220/237046.html}{\textcolor{blue}{Nuclear Science and Technology Youth Forum}}
\newline{} 22 December 2023 (invited talk)}

\cvitem{6}{\textit{Real-time dynamics of the $O(N)$ scalar theory within the fRG approach}\newline{} \href{https://sites.google.com/view/funqcd22/}{\textcolor{blue}{FunQCD: from first principles to effective theories}}\newline{} 13 to 17 June 2022 (online)}



\cvitem{7}{\textit{Real time dynamics of the $O(4)$ scalar theory within the fRG approach}\newline{} \href{https://indico.ihep.ac.cn/event/13478/}{\textcolor{blue}{The 14th Workshop on QCD Phase Transition and Relativistic Heavy-ion Physics (QPT 2021)}}\newline{}26 to 30 July 2021}

\cvitem{8}{\textit{Real time dynamics of the $O(4)$ scalar theory within the fRG approach}\newline{} \href{https://funqcd.blogs.uv.es/}{\textcolor{blue}{FunQCD: from first principles to effective theories}}\newline{}29 March to 1 April 2021 (Flash talk)}


%-----------TECHNICAL SKILLS-----------
\section{Technical Skills}
\cvitem{}{GitHub Profile: \href{https://github.com/Yangyang-Tan}{\textcolor{blue}{https://github.com/Yangyang-Tan}}}
\cvitem{Wolfram Language \& Mathematica}{Expert -- 12 years of experience}
\cvitem{Julia}{Proficient -- 5 years of experience}
\cvitem{Python(PyTorch)}{Proficient -- 2 years of experience}
\cvitem{Large-scale parallel computing and GPU computing (based on Julia)}{Proficient -- 5 years of experience}

%-----------PhD SUPERVISOR-----------
\section{PhD Supervisor}
\cvitem{}{Wei-jie Fu, Professor of School of Physics, Dalian University of Technology}
\cvitem{Address}{No.2 Linggong Road, Ganjingzi District, Dalian City, Liaoning Province, P.R.C., 116024}
\cvitem{Email}{\href{mailto:wjfu@dlut.edu.cn}{wjfu@dlut.edu.cn}}
\cvitem{Telephone}{(86)13478905891}

\end{document}